\documentclass[UTF8]{ctexart}
\usepackage{xeCJK}
\usepackage{geometry}
\geometry{left=2cm,right=2cm,top=2.5cm,bottom=2.5cm}
\setlength{\headheight}{12.7pt}

\usepackage{titlesec}
\titleformat{\section}
  {\normalfont\large\bfseries}
  {\thesection}
  {1em}
  {}
\titlespacing*{\section}
  {0pt}
  {3.5ex plus 1ex minus .2ex}
  {2.3ex plus .2ex}

\usepackage{amsmath}
\usepackage{amssymb}
\usepackage{amsthm}
\usepackage{enumitem}
\setlist[enumerate]{itemsep=0pt,topsep=0pt,parsep=0pt,leftmargin=0pt,itemindent=2.5em}
\setlist[itemize]{itemsep=0pt,topsep=0pt,parsep=0pt,leftmargin=0pt,itemindent=2.5em}
\usepackage{fancyhdr}
\pagestyle{fancy}
\fancyhf{}
\usepackage{graphicx}
\usepackage{hyperref}
\usepackage{indentfirst}
\usepackage{lmodern}


\begin{document}
\setlength{\abovedisplayskip}{1pt}
\setlength{\belowdisplayskip}{1pt}

\section{良序关系}

\hypertarget{sec:定义1.1}{}
\noindent\textbf{定义1.1 (良序关系)}

给定集合$a$和$a$上的二元关系$\leqslant$. 如果以下五条成立:
\begin{enumerate}
    \item 自反性: $\forall x\in a:\,x\leqslant x$,
    \item 传递性: $\forall x,y,z\in a:\,
    x\leqslant y\land y\leqslant z\to x\leqslant z$,
    \item 反对称性: $\forall x,y\in a:\,x\leqslant y\land y\leqslant x\to x=y$,
    \item 可比性: $\forall x,y\in a:\,x\leqslant y\lor y\leqslant x$,
    \item 良基性: 对$a$的任意非空子集$b$, 存在$m\in b$使得
    $\forall x\in b:\,m\leqslant x$.
\end{enumerate}
那么称$\leqslant$是$a$上的一个\textbf{良序关系}.
也称$(a,\leqslant)$是一个\textbf{良序结构}.

\vspace{10pt}

给定良序结构$(a,\leqslant)$, 它衍生的$<$满足以下四条, 由此称之为\textbf{严格良序}:
\begin{enumerate}
    \item 反自反性: $\forall x\in a:\,x\nless x$,
    \item 传递性: $\forall x,y,z\in a:\,x<y\land y<z\to x<z$,
    \item 三歧性: $\forall x,y\in a:\,x<y\lor x=y\lor x>y$,
    \item 良基性: 对$a$的任意非空子集$b$, 存在$m\in b$使得
    $\forall x\in b:\,x\nless m$.
\end{enumerate}
因为良序没有要求子集存在最大元, 所以此时$\geqslant$不一定是良序, $>$也不一定是严格良序.

\vspace{10pt}

\hypertarget{sec:定理1.1}{}
\noindent\textbf{定理1.1}

有限集上的全序一定是良序.

\noindent{\kaishu 证明.}

有限集上的全序的非空子集一定有最小元, 所以是良序. \hfill $\qedsymbol$

\vspace{10pt}

\hypertarget{sec:定理1.2}{}
\noindent\textbf{定理1.2 (无穷降链)}

给定良序结构$(a,\leqslant)$和$a$中的序列$(f_n)_{n\in\omega}$. 如果它是递降序列, 即
\[
    \forall n\in\omega:f_{n+1}\leqslant f_n,
\]
那么存在自然数$N$使得$\forall n\geqslant N:f_n=f_N$.

\noindent{\kaishu 证明.}

记$\{f_n\mid n\in\omega\}$的最小元为$m$, 存在自然数$N$使得$f_N=m$.
然后对$n\geqslant N$归纳证明$f_n=f_N$.

首先$f_N=f_N$; 其次任取自然数$n\geqslant N$, 假设$f_n=f_N$, 那么
\[
    f_{n+1}\leqslant f_n\quad\land\quad f_N\leqslant f_{n+1},
\]
从而$f_{n+1}=f_N$. \hfill $\qedsymbol$

\vspace{10pt}

\hypertarget{sec:定理1.3}{}
\noindent\textbf{定理1.3}

给定良序结构$(a,\leqslant)$和真前段$s$, 存在$t\in a$使得$s=a_t$.

\noindent{\kaishu 证明.}

$a\setminus s$非空, 记$t$为其最小元.

\begin{enumerate}
    \item 对任意的$x\in s$, 假设$t\leqslant x$,
    那么由$s$是前段得$t\in s$, 矛盾. 所以$x<t$, 即$x\in a_t$.
    \item 对任意的$x\in a_t$, 假设$x\in a\setminus s$则同时有$x<t$和$t\leqslant x$,
    矛盾. 从而$x\in s$.
\end{enumerate}
综上即得$s=a_t$. \hfill $\qedsymbol$





\newpage

\section{良序的归纳}

\hypertarget{sec:定理2.1}{}
\noindent\textbf{定理2.1 (良基归纳法)}

给定$a$和其上的良基关系$\leqslant$和命题$p(x)$. 如果
\[
    \forall x\in a:\quad(\,\forall y<x:p(y)\,)\to p(x),
\]
那么对任意的$x\in a$, $p(x)$成立.

\noindent{\kaishu 证明.}

假设$b=\{x\in a\mid\neg\,p(x)\}$非空, 记极小元为$m$.

此时对任意的$y<m$有$y\notin b$, 即$p(y)$成立. 从而$p(m)$成立, 与$m\in b$矛盾.
\hfill $\qedsymbol$

\vspace{10pt}

下面先使用超限归纳法证明一些良序同构相关的定理.

\vspace{10pt}

\hypertarget{sec:定理2.2}{}
\noindent\textbf{定理2.2 (良序同构唯一)}

给定良序结构$(a,\leqslant_a)$和$(b,\leqslant_b)$.
如果$f,g:(a,\leqslant_a)\cong(b,\leqslant_b)$, 那么$f=g$.

\noindent{\kaishu 证明.}

对$x\in a$归纳证明$f(x)=g(x)$. 任取$x\in a$, 假设$\forall y<_ax:f(y)=g(y)$.

再假设$f(x)<_bg(x)$, 记$y=g^{-1}(f(x))$, 则
\[
    y<_ag^{-1}(g(x))=x\implies f(y)=g(y)=f(x)\implies y=x,
\]
矛盾. 同理假设$f(x)>_bg(x)$也矛盾, 所以$f(x)=g(x)$. \hfill $\qedsymbol$

\vspace{10pt}

\hypertarget{sec:定理2.3}{}
\noindent\textbf{定理2.3}

给定良序结构$(a,\leqslant)$和前段$s,s'$.
如果$(s,\leqslant)\cong(s',\leqslant)$, 那么$s=s'$.

\noindent{\kaishu 证明.}

取$f:(s,\leqslant)\cong(s',\leqslant)$, 对$x\in s$归纳证明$f(x)=x$.
任取$x\in s$, 假设$\forall y<x:f(y)=y$.

\begin{enumerate}
    \item 假设$f(x)<x$, 那么依假设有$f(f(x))=f(x)$, 从而$f(x)=x$, 矛盾.
    \item 假设$x<f(x)$, 那么$x\in s'$, 从而$f^{-1}(x)<x$,
    依假设有$x=f(f^{-1}(x))=f^{-1}(x)$, 矛盾.
\end{enumerate}
综上$f(x)=x$, 归纳得$\forall x\in s:f(x)=x$. 所以$s=s'$. \hfill $\qedsymbol$





\newpage

最后给出良序集基本定理, 它说明任意两个良序都是可比的.

\vspace{10pt}

\hypertarget{sec:定理2.4}{}
\noindent\textbf{定理2.4 (良序集基本定理)}

任给良序结构$(a,\leqslant_a)$和$(b,\leqslant_b)$, 下面三条有且只有一条成立:
\begin{enumerate}
    \item $(a,\leqslant_a)\cong(b,\leqslant_b)$.
    \item 存在$x\in a$使得$(a_x,\leqslant_a)\cong(b,\leqslant_b)$.
    \item 存在$y\in b$使得$(a,\leqslant_a)\cong(b_y,\leqslant_b)$.
\end{enumerate}

\vspace{3pt}

\noindent{\kaishu 证明. 下面写序同构时, 只写集合, 省略关系符号.}

\vspace{3pt}

定义从$a$到$b$的二元关系
\[
    r=\{(x,y)\in a\times b\mid a_x\cong b_y\}.
\]

首先证明$r$是双射. 如果$(x,y)\in r$且$(x,y')\in r$, 那么存在
\[
    f:a_x\cong b_y,\quad g:a_x\cong b_{y'},
\]
于是
\[
    g\circ f^{-1}:b_y\cong b_{y'}\quad\implies\quad b_y=b_{y'}\implies y=y'.
\]
同理如果$(x,y)\in r$且$(x',y)\in r$, 那么$x=x'$. 所以$r$是双射.

\vspace{3pt}

然后证明$r$是序同构. 对任意的$x\in\mathrm{Dom}\,r$和$x'<_ax$,
取$f:a_x\cong b_{rx}$, 从而
\[
    (a_x)_{x'}\cong(b_{rx})_{fx'}\implies a_{x'}\cong b_{fx'}
    \implies(x',f(x'))\in r.
\]
所以$r(x')=f(x')\in b_{rx}$, 即$r(x')<_br(x)$.

另外上证了$x'\in\mathrm{Dom}\,r$, 所以$\mathrm{Dom}\,r$是$a$的前段.
同理$\mathrm{Ran}\,r$是$b$的前段.

\vspace{3pt}

最后, 下面四条有且只有一条成立:
\begin{enumerate}
    \item $\mathrm{Dom}\,r=a,\mathrm{Ran}\,r=b$, 从而$a\cong b$.
    \item $\mathrm{Dom}\,r\neq a,\mathrm{Ran}\,r=b$,
    则存在$x\in a$使得$\mathrm{Dom}\,r=a_x$, 从而$a_x\cong b$.
    \item $\mathrm{Dom}\,r=a,\mathrm{Ran}\,r\neq b$,
    则存在$y\in b$使得$\mathrm{Ran}\,r=b_y$, 从而$a\cong b_y$.
    \item $\mathrm{Dom}\,r\neq a,\mathrm{Ran}\,r\neq b$,
    则存在$x\in a$和$y\in b$使得$\mathrm{Dom}\,r=a_x,\mathrm{Ran}\,r=b_y$, 此时
    \[
        a_x\cong b_y\implies(x,y)\in r\implies x\in a_x\land y\in b_y,
    \]
    矛盾. \hfill $\qedsymbol$
\end{enumerate}

\end{document}